\documentclass[12pt,a4paper]{article}
\usepackage[utf8]{inputenc}
\usepackage[T1]{fontenc}
\usepackage[spanish]{babel}
\usepackage{geometry}
\usepackage{xcolor}
\usepackage{enumitem}

\geometry{a4paper, margin=1in}

\title{Informe de Estado de Pull Requests \\ \large Curso: Seguridad Informática}
\author{Prof. César Rodríguez}
\date{\today}

\begin{document}

\maketitle

\section*{Resumen de Integración}
Este documento presenta el estado actual de los \textit{Pull Requests} enviados por los distintos grupos de desarrollo de la Suite de Auditoría. Solo se listan los grupos cuyos módulos han sido evaluados, corregidos y fusionados con éxito en la rama principal, cumpliendo con el contrato de datos establecido.

\section*{Grupos con Pull Requests Aprobados}

\subsection*{Grupo 2 - OSINT (\texttt{osint.py})}
\textbf{Estado:} \textcolor{green!70!black}{\textbf{Aprobado e Integrado}} \\
\textbf{Calificación:} \textbf{10/10} (Entrega correcta desde el principio)
\begin{itemize}[label=$\bullet$]
    \item \textbf{Estudiante 1 (WHOIS):} Implementó la consulta WHOIS para extraer fechas de registro, nombre de la organización y servidores DNS.
    \item \textbf{Estudiante 2 (Dorks - Subdominios):} Desarrolló la automatización de Google Dorks para descubrir subdominios omitiendo el subdominio principal.
    \item \textbf{Estudiante 3 (Dorks - Archivos Expuestos):} Implementó la búsqueda de archivos de configuración (\texttt{.env}, \texttt{.ini}) gestionando correctamente el límite de peticiones (Rate Limit HTTP 429).
\end{itemize}

\subsection*{Grupo 3 - Discovery (\texttt{discovery.py})}
\textbf{Estado:} \textcolor{green!70!black}{\textbf{Aprobado e Integrado}} \\
\textbf{Calificación:} \textbf{8/10} (Penalización por errores cometidos por el Estudiante 2)
\begin{itemize}[label=$\bullet$]
    \item \textbf{Estudiante 1 (Ping Sweep):} Implementación correcta de escaneo ICMP, incluyendo la actualización para soporte multiplataforma (Windows y Linux/macOS).
    \item \textbf{Estudiante 2 (TCP SYN Ping):} Lógica de escaneo con Scapy funcional. Se destaca la excelente práctica de enviar paquetes RST para cerrar conexiones \textit{Half-Open}. El contrato de datos fue corregido y verificado.
    \item \textbf{Estudiante 3 (TCP ACK Ping):} Excelente implementación del escaneo ACK detectando las respuestas con bandera RST (\texttt{0x4}).
\end{itemize}
\textit{Nota:} El código se fusionó en la rama principal tras aplicar correcciones menores de sintaxis e importación de dependencias por parte del revisor.

\section*{Grupos Pendientes de Revisión}
\begin{itemize}[label=$\bullet$]
    \item \textbf{Grupo 4 - Scanning (\texttt{scanning.py}):} A la espera de Pull Request unificado.
    \item \textbf{Grupo 1 - Reconocimiento DNS (\texttt{dns\_recon.py}):} A la espera de Pull Request unificado.
\end{itemize}

\vspace{1cm}
\hrule
\vspace{0.5cm}
\noindent \textit{El presente informe se actualizará conforme los grupos restantes finalicen sus entregas y superen la validación estricta de \texttt{auditoria.py}.}

\end{document}